% Part: sets-functions-relations
% Chapter: functions
% Section: inverses
% Hindi translation (hi-Deva-IN) of Open Logic commit 9620cc73f9c8e0ad003c514a5d3748f29611c4c0.

\documentclass[../../../include/open-logic-section]{subfiles}

\begin{document}

\olfileid[hi]{sfr}{fun}{inv}
\olsection{फलनों के प्रतिलोम}

\begin{explain}
हम फलनों को प्रतिचित्रण मानते हैं। तब फलनों के बारे में एक स्वाभाविक प्रश्न
यह है कि क्या प्रतिचित्रण को “उलटा” जा सकता है। उदाहरणार्थ, उत्तराधिकारी फलन
$f(x) = x + 1$ को इस अर्थ में उलटा जा सकता है कि फलन $g(y) = y - 1$,
$f$ के किए को “पूर्ववत्” कर देता है।

पर हमें सावधान रहना होगा। यद्यपि~$g$ की परिभाषा एक फलन
$\Int \to \Int$ परिभाषित करती है, वह कोई \emph{फलन} $\Nat
\to \Nat$ परिभाषित नहीं करती, क्योंकि $g(0) \notin \Nat$। इसलिए सरल
स्थितियों में भी यह एकदम स्पष्ट नहीं होता कि किसी फलन को उलटा जा सकता है या
नहीं; यह उसके प्रांत और सहप्रांत पर निर्भर हो सकता है।

फलन के प्रतिलोम की धारणा इस बात को अधिक सटीक बनाती है।
\end{explain}

\begin{defn}
फलन $g \colon B \to A$, फलन $f
\colon A \to B$ का \emph{प्रतिलोम} है, यदि $f(g(y)) = y$ और
$g(f(x)) = x$, सभी $x \in A$ और $y \in B$ के लिए हों।
\end{defn}

यदि $f$ का प्रतिलोम~$g$ हो, तो हम प्रायः $f^{-1}$ को~$g$ के स्थान पर लिखते हैं।

\begin{explain}
अब हम निर्धारित करेंगे कि किन फलनों के प्रतिलोम होते हैं। $f\colon A \to B$
के प्रतिलोम का एक स्वाभाविक उम्मीदवार $g\colon B \to A$ है, जिसे इस प्रकार
“परिभाषित” किया जाए:
\[
g(y) = \text{वह “एकमात्र” $x$ जिसके लिए $f(x) = y$ है।}
\]
पर “परिभाषित” (और “एकमात्र”) के चारों ओर लगे उद्धरण-चिह्न बताते हैं कि यह
वास्तव में परिभाषा नहीं है। कम-से-कम पूर्ण व्यापकता में यह सदा काम नहीं
करेगी। इस कथन से कोई फलन निर्धारित हो, इसके लिए ऐसा ठीक एक~$x$ होना चाहिए
जिसके लिए $f(x) =
y$---अर्थात~$g$ का निर्गत अद्वितीय रूप से निर्धारित होना
चाहिए। इसके अतिरिक्त, उसे प्रत्येक $y \in B$ के लिए निर्धारित होना चाहिए।
यदि $x_1$ और $x_2 \in
A$ ऐसे हों कि $x_1 \neq x_2$, किंतु
$f(x_1) = f(x_2)$, तो $g(y)$ उस $y = f(x_1) = f(x_2)$ के लिए अद्वितीय
रूप से निर्धारित नहीं होगा। और यदि ऐसा कोई~$x$ हो ही नहीं जिसके लिए $f(x) = y$,
तो $g(y)$ बिल्कुल निर्धारित नहीं होगा। दूसरे शब्दों में, $g$ के परिभाषित
होने के लिए $f$ का !!{injective} और !!{surjective} दोनों होना आवश्यक है।

आइए इसे क्रमशः समझें। प्रश्न को दो भागों में बाँटेंगे। दिए गए फलन
$f\colon A \to B$ के लिए ऐसा फलन $g\colon B \to A$ कब होता है कि
$g(f(x)) = x$? ऐसा $g$, $f$ के किए को “पूर्ववत्” करता है और उसे~$f$ का
\emph{वाम प्रतिलोम} कहते हैं। दूसरा, ऐसा फलन $h\colon B \to A$ कब होता है
कि $f(h(y)) = y$? ऐसे $h$ को~$f$ का \emph{दक्षिण प्रतिलोम} कहते हैं---इस
बार $f$, $h$ के किए को “पूर्ववत्” करता है।
\end{explain}

\begin{prop}
यदि $f\colon A \to B$ !!{injective} है, तो कोई फलन
$g\colon B \to A$ ऐसा होता है जो~$f$ का \emph{वाम प्रतिलोम} है और जिसके लिए
$g(f(x)) = x$, सभी $x
\in A$ के लिए।
\end{prop}

\begin{proof}
मान लीजिए कि $f\colon A \to B$ !!{injective} है। किसी $y \in B$ पर विचार
कीजिए। यदि $y \in \ran{f}$, तो कोई $x \in A$ ऐसा है कि $f(x) = y$। चूँकि
$f$ !!{injective} है, ऐसा केवल एक~$x \in A$ है। तब हम परिभाषित कर सकते हैं:
$g(y) = x$, अर्थात $g(y)$ वह “एकमात्र” $x \in A$ है जिसके लिए $f(x)
= y$। यदि $y \notin \ran{f}$, तो हम उसे किसी भी~$a \in A$ पर प्रतिचित्रित
कर सकते हैं। अतः कोई $a \in A$ चुनकर $g\colon B \to A$ को इस प्रकार
परिभाषित कर सकते हैं:
\[
g(y) = \begin{cases}
    x & \text{यदि $f(x) = y$}\\
    a & \text{यदि $y \notin \ran{f}$।}
\end{cases}
\]
यह सभी $y \in B$ के लिए परिभाषित है, क्योंकि ऐसे प्रत्येक
$y \in \ran{f}$ के लिए ठीक एक $x \in A$ है जिसके लिए $f(x) = y$। परिभाषा
से, यदि $y = f(x)$, तो $g(y) = x$, अर्थात $g(f(x)) = x$।
\end{proof}

\begin{prob}
दिखाइए कि यदि $f\colon A \to B$ का वाम प्रतिलोम~$g$ है, तो $f$
!!{injective} है।
\end{prob}

\begin{prop}
    यदि $f\colon A \to B$ !!{surjective} है, तो कोई फलन
    $h\colon B \to A$ ऐसा होता है जो~$f$ का \emph{दक्षिण प्रतिलोम} है और जिसके लिए
    $f(h(y)) =
    y$ सभी~$y \in B$ के लिए।
\end{prop}

\begin{proof}
मान लीजिए कि $f\colon A \to B$ !!{surjective} है। किसी $y \in
B$ पर विचार कीजिए। चूँकि $f$ !!{surjective} है, कोई $x_y \in A$ ऐसा है कि
$f(x_y)
= y$। तब हम $h(y) = x_y$ परिभाषित कर सकते हैं; अर्थात प्रत्येक
$y \in B$ के लिए कोई ऐसा $x \in A$ चुनते हैं कि $f(x) = y$। चूँकि $f$
!!{surjective} है, चुनने के लिए कम-से-कम एक ऐसा अवयव सदा होता है।\footnote{चूँकि
$f$ !!{surjective} है, प्रत्येक~$y \in B$ के लिए समुच्चय
$\Setabs{x}{f(x) = y}$ अरिक्त है।~$h$ की हमारी परिभाषा में इन समुच्चयों में
से प्रत्येक से एक $x$ चुनना पड़ता है। ऐसा चयन सदा संभव है, यह वास्तव में
स्पष्ट नहीं है---इन चयनों को कर पाने की संभावना को बस एक अभिगृहीत मान लिया
जाता है। दूसरे शब्दों में, यह प्रतिज्ञप्ति तथाकथित चयन के अभिगृहीत को मानती
है; इस विषय पर हम \oliflabeldef{sth:choice::chap}{\olref[sth][choice][]{chap}
में फिर विचार करेंगे}{अभी विस्तार से विचार नहीं करेंगे}। किंतु अनेक विशिष्ट
स्थितियों में, जैसे $A = \Nat$ हो या परिमित हो, अथवा $f$ !!{bijective} हो,
तब चयन के अभिगृहीत की आवश्यकता नहीं होती। (विशेषतः जब $f$ !!{bijective} हो,
तब प्रत्येक $y \in B$ के लिए समुच्चय $\Setabs{x}{f(x) = y}$ में ठीक एक
!!{element} होता है, इसलिए कोई चयन करना ही नहीं पड़ता।)}
परिभाषा से, यदि $x = h(y)$, तो $f(x) = y$; अर्थात किसी भी $y \in B$ के
लिए $f(h(y)) = y$।
\end{proof}

\begin{prob}
दिखाइए कि यदि $f\colon A \to B$ का दक्षिण प्रतिलोम~$h$ है, तो $f$
!!{surjective} है।
\end{prob}

\begin{explain}
  पिछले प्रमाण के विचारों को मिलाने पर अब हमें मिलता है कि प्रत्येक
  !!{bijection} का एक प्रतिलोम होता है; अर्थात एक ही फलन~$f$ का वाम और
  दक्षिण, दोनों प्रतिलोम होता है।
\end{explain}

\begin{prop}\ollabel{prop:bijection-inverse}
यदि $f\colon A \to B$ !!{bijective} है, तो ऐसा फलन
$f^{-1}\colon B \to A$ होता है कि सभी $x \in A$ के लिए
$f^{-1}(f(x)) = x$ और सभी $y \in B$ के लिए $f(f^{-1}(y)) = y$।
\end{prop}

\begin{proof}
अभ्यास।
\end{proof}

\begin{prob}
\olref[sfr][fun][inv]{prop:bijection-inverse} सिद्ध कीजिए। आपको~$f^{-1}$
परिभाषित करना है, यह दिखाना है कि वह एक फलन है, और यह दिखाना है कि वह~$f$
का प्रतिलोम है; अर्थात $f^{-1}(f(x)) = x$ और $f(f^{-1}(y)) = y$, सभी
$x \in A$ और $y \in B$ के लिए।
\end{prob}

\begin{explain}
प्रतिलोम प्राप्त करने की एक थोड़ी अधिक व्यापक विधि भी है। हमने
\olref[kin]{sec} में देखा कि प्रत्येक फलन $f$, !!a{surjection} $f'
\colon A \to \ran{f}$ इस प्रकार प्रेरित करता है कि $f'(x) = f(x)$ सभी
$x \in A$ के लिए। स्पष्टतः, यदि $f$ !!{injective} है,
तो $f'$ !!{bijective} है; अतः \olref{prop:bijection-inverse} के अनुसार उसका
एक अद्वितीय प्रतिलोम है। संकेतन का बहुत मामूली दुरुपयोग करते हुए हम कभी-कभी
$f'$ के प्रतिलोम को बस “$f$ का प्रतिलोम” कहते हैं।
\end{explain}

\begin{prop}\ollabel{prop:left-right}%
  दिखाइए कि यदि $f\colon A \to B$ का वाम प्रतिलोम~$g$ और दक्षिण
  प्रतिलोम~$h$ है, तो $h = g$।
\end{prop}

\begin{proof}
  अभ्यास।
\end{proof}

\begin{prob}
  \olref[sfr][fun][inv]{prop:left-right} सिद्ध कीजिए।
\end{prob}

\begin{prop}\ollabel{prop:inverse-unique}
प्रत्येक फलन~$f$ का अधिक-से-अधिक एक प्रतिलोम होता है।
\end{prop}

\begin{proof}
  मान लीजिए कि $g$ और $h$, दोनों~$f$ के प्रतिलोम हैं। तब विशेषतः $g$,
  $f$ का वाम प्रतिलोम और $h$, उसका दक्षिण प्रतिलोम है।
  \olref{prop:left-right} से $g = h$।
\end{proof}

\end{document}
